

















\section{Details on Differentiable Transplant Operators}
\label{app:designers}

This section gives the implementation details for the differentiable shared-basis transplant operators used in the cross-operator validation: FOCUS~\citep{dobler2023focus}, CLP~\citep{ostendorff2023efficient}, and WECHSEL~\citep{minixhofer2021wechsel}.
Each operator computes donor-side weights and synthesizes the base row as a weighted combination of base anchors. This final coefficient-reuse step is the shared structure analyzed in Appendix~\ref{app:nonlinear-shared-anchor}.
In cases where the operator involves a linear transformation (e.g., affine alignment or projection), we denote the corresponding transformation matrix as $\mathbf{A}$.
Differentiability lets us optimize the donor embedding $\mathbf{x}_d$ directly via backpropagation, rather than using the discrete support search used by OMP.

\noindent\textbf{Optimization Objective.}
Consistent with the general framework (Eq.~\ref{eq:designer_obj}), we set the learnable parameters to be the raw donor embedding itself ($\theta = \mathbf{x}_d$). We minimize the following loss:
\begin{equation}
\label{eq:diff_obj}
    \mathcal{L}(\mathbf{x}_d)
    \;=\;
    \underbrace{\|f(\mathbf{x}_d) - \boldsymbol{\mu}_{\text{base}}\|_2^2}_{\text{Base Salience}}
    \;+\;
    \lambda \underbrace{\|U \mathbf{x}_d\|_2^2}_{\text{Donor Inertness}}
    \;+\;
    \rho \underbrace{(\|\mathbf{x}_d\|_2-\nu_d)^2}_{\text{Target-norm regularization}},
\end{equation}
where $\rho>0$ controls norm preservation and $\nu_d$ is the median natural donor embedding norm for the corresponding model.
We solve Eq.~\ref{eq:diff_obj} using a unified gradient-based approach (Adam) for all operators described below.
\subsection{CLP: Contrastive Linear Projection}
\label{sec:clp-breaker}

\noindent\textbf{Transplant mechanism.}
For the barycentric CLP variant~\citep{ostendorff2023efficient}, weights are computed by rectifying cosine similarities and renormalizing.
Using normalized donor input $\widehat{\mathbf{x}}_d$ and anchors $\widehat{\boldsymbol{\phi}}^{(d)}_j$, the weights are:
\begin{equation}
\label{eq:clp_w}
w_j(\mathbf{x}_d)
\;=\;
\frac{\mathrm{ReLU}\!\big(\langle \widehat{\mathbf{x}}_d, \widehat{\boldsymbol{\phi}}^{(d)}_j\rangle\big)}
{\epsilon + \sum_{k\in\mathcal{T}} \mathrm{ReLU}\!\big(\langle \widehat{\mathbf{x}}_d, \widehat{\boldsymbol{\phi}}^{(d)}_k\rangle\big)}.
\end{equation}
The transplanted embedding is synthesized as:
\begin{equation}
\label{eq:clp_func}
f_{\text{CLP}}(\mathbf{x}_d)
\;=\;
\sum_{j\in\mathcal{T}} w_j(\mathbf{x}_d)\,\boldsymbol{\phi}^{(b)}_j.
\end{equation}

\noindent\textbf{Designer optimization.}
We minimize the differentiable objective (Eq.~\ref{eq:diff_obj}) with $f=f_{\text{CLP}}$.
Because anchors with negative cosine similarity yield zero gradient through the ReLU function, we employ a \textbf{non-negative reverse projection} for initialization to prevent vanishing gradients.
We choose a small anchor set $\mathcal{S}$ near $\boldsymbol{\mu}_{\text{base}}$ in the base dictionary and solve
\[
\boldsymbol{\alpha}^\star
=
\arg\min_{\boldsymbol{\alpha}\ge 0}
\left\|\sum_{j\in\mathcal{S}}\alpha_j\boldsymbol{\phi}^{(b)}_j-\boldsymbol{\mu}_{\text{base}}\right\|_2^2.
\]
We then set $\mathbf{x}_{\text{init}}=\sum_{j\in\mathcal{S}}\alpha^\star_j\boldsymbol{\phi}^{(d)}_j$ so the starting vector positively activates the target anchors.

\subsection{WECHSEL: Alignment and Neighbors}
\label{sec:wechsel-breaker}

\noindent\textbf{Transplant mechanism.}
WECHSEL~\citep{minixhofer2021wechsel} combines a global affine alignment with a local $k$-NN mixture over the \emph{base} vocabulary.
Let $\boldsymbol{\mu}_{\text{src}}, \boldsymbol{\mu}_{\text{tgt}}$ be the centroids of the donor and base embeddings, respectively, and let $\mathbf{A}_{\mathrm{WEC}}$ be the alignment map from donor to base coordinates learned from shared tokens.
We define the \emph{aligned proxy} of a donor input $\mathbf{x}_d$ as:
\begin{equation}
\label{eq:wechsel_align}
\tilde{\mathbf{x}}(\mathbf{x}_d) \;=\; \mathbf{A}_{\mathrm{WEC}}(\mathbf{x}_d - \boldsymbol{\mu}_{\text{src}}) + \boldsymbol{\mu}_{\text{tgt}}.
\end{equation}
Let $\mathcal{N}_k(\tilde{\mathbf{x}})$ be the set of indices of the $k$ base anchors $\{\boldsymbol{\phi}^{(b)}_v\}$ nearest to $\tilde{\mathbf{x}}$ (under the specific distance metric used by WECHSEL), and let $w_v(\tilde{\mathbf{x}})$ be the corresponding softmax weights.
The transplanted embedding is synthesized as:
\begin{equation}
\label{eq:wechsel_func}
f_{\text{WEC}}(\mathbf{x}_d)
\;=\;
\sum_{v\in\mathcal{N}_k(\tilde{\mathbf{x}})} w_v(\tilde{\mathbf{x}}(\mathbf{x}_d))\,\boldsymbol{\phi}^{(b)}_v.
\end{equation}

\noindent\textbf{Designer optimization.}
Because the synthesis basis $\mathcal{N}_k(\tilde{\mathbf{x}})$ changes discretely with $\mathbf{x}_d$, the objective is non-convex and piecewise-smooth.
We initialize in the aligned base space and invert the affine map. First choose a small anchor set $\mathcal{S}_{\mathrm{WEC}}$ near $\boldsymbol{\mu}_{\text{base}}$ and solve
\[
\boldsymbol{\alpha}^\star
=
\arg\min_{\boldsymbol{\alpha}\in\Delta^{|\mathcal{S}_{\mathrm{WEC}}|}}
\left\|\sum_{j\in\mathcal{S}_{\mathrm{WEC}}}\alpha_j\boldsymbol{\phi}^{(b)}_j-\boldsymbol{\mu}_{\text{base}}\right\|_2^2.
\]
Let $\mathbf{z}_{\text{init}}=\sum_{j\in\mathcal{S}_{\mathrm{WEC}}}\alpha_j^\star\boldsymbol{\phi}^{(b)}_j$. We set
\[
\mathbf{x}_{\text{init}}
=
\boldsymbol{\mu}_{\text{src}}+\mathbf{A}_{\mathrm{WEC}}^\dagger(\mathbf{z}_{\text{init}}-\boldsymbol{\mu}_{\text{tgt}}),
\]
then apply the donor preconditioner $\mathbf{x}_{\text{init}}\leftarrow(I+\lambda U^\top U)^{-1}\mathbf{x}_{\text{init}}$ before Adam refinement.
We minimize the unified differentiable objective (Eq.~\ref{eq:diff_obj}) by substituting the operator $f_{\text{WEC}}$:
\begin{equation}
\label{eq:wechsel_obj}
\min_{\mathbf{x}_d} \; \mathcal{L}(\mathbf{x}_d)
\;=\;
 \big\|f_{\text{WEC}}(\mathbf{x}_d) - \boldsymbol{\mu}_{\text{base}}\big\|_2^2
\;+\;
\lambda \|U \mathbf{x}_d\|_2^2
\;+\;
\rho(\|\mathbf{x}_d\|_2-\nu_d)^2.
\end{equation}
The alignment transformation in Eq.~\ref{eq:wechsel_align} is embedded within $f_{\text{WEC}}$, so refinement optimizes the raw donor vector $\mathbf{x}_d$ directly.

\subsection{FOCUS: Soft Anchor Interpolation}
\label{sec:focus-breaker}

\noindent\textbf{Transplant mechanism.}
FOCUS~\citep{dobler2023focus} reconstructs tokens via similarity-weighted interpolation of shared semantic anchors.
Let $\mathcal{T}$ be the set of shared tokens. We denote the donor and base anchor matrices as $\Phi_d$ and $\Phi_b$, where the $j$-th row corresponds to the shared token $j \in \mathcal{T}$.
For numerical stability, the mechanism operates on normalized donor anchors $\widehat{\boldsymbol{\phi}}^{(d)}_j = \boldsymbol{\phi}^{(d)}_j / \|\boldsymbol{\phi}^{(d)}_j\|_2$.
Given a donor input $\mathbf{x}_d$, the operator computes attention weights using a softmax kernel with inverse temperature $\beta_{\mathrm{sm}}$:
\begin{equation}
\label{eq:focus_w}
w_j(\mathbf{x}_d)
\;=\;
\frac{\exp\!\big(\beta_{\mathrm{sm}}\,\langle \widehat{\mathbf{x}}_d, \widehat{\boldsymbol{\phi}}^{(d)}_j\rangle\big)}
{\sum_{k\in\mathcal{T}} \exp\!\big(\beta_{\mathrm{sm}}\,\langle \widehat{\mathbf{x}}_d, \widehat{\boldsymbol{\phi}}^{(d)}_k\rangle\big)},
\end{equation}
where $\widehat{\mathbf{x}}_d$ is the normalized input. The transplanted embedding is then synthesized as the weighted sum of base anchors:
\begin{equation}
\label{eq:focus_func}
f_{\text{FOC}}(\mathbf{x}_d)
\;=\;
\sum_{j\in\mathcal{T}} w_j(\mathbf{x}_d)\,\boldsymbol{\phi}^{(b)}_j.
\end{equation}

\noindent\textbf{Designer optimization.}
Due to the nonlinearity of the attention weights $w_j(\mathbf{x}_d)$, the objective is non-convex.
We minimize the unified differentiable objective (Eq.~\ref{eq:diff_obj}) by substituting $f=f_{\text{FOC}}$ and employing gradient-based optimization.

To avoid local minima, we employ a \textbf{reverse-projection initialization}. We identify a small subset of base anchors $\mathcal{S}\subset\mathcal{T}$ (e.g., the $k$-nearest neighbors of $\boldsymbol{\mu}_{\text{base}}$ in $\Phi_b$) and solve for the optimal reconstruction coefficients in the base space:
\begin{equation}
\label{eq:focus_init_ls}
\boldsymbol{\alpha}^\star
\;=\;
\arg\min_{\boldsymbol{\alpha}\in\Delta^{|\mathcal{S}|}}
\Big\|\sum_{j\in\mathcal{S}} \alpha_j \boldsymbol{\phi}^{(b)}_j - \boldsymbol{\mu}_{\text{base}}\Big\|_2^2.
\end{equation}
We then initialize the donor vector by mapping these coefficients back to the donor space: $\mathbf{x}_{\text{init}} = \sum_{j\in\mathcal{S}} \alpha^\star_j \boldsymbol{\phi}^{(d)}_j$.
Finally, we apply the donor suppression constraint as a one-shot preconditioner before optimization: $\mathbf{x}_{\text{init}} \leftarrow (I+\lambda U^\top U)^{-1}\mathbf{x}_{\text{init}}$.
\subsection{SER results on differentiable transplant operators}
\label{sec:altops-results}

We now evaluate whether the breaker-token designer transfers to the three shared-basis operator variants described above. Table~\ref{tab:app:altops:ser} reports the cross-operator validation rows for CLP, WECHSEL, and FOCUS, including the evaluated $\lambda$ values and SER for both the attacked base and the patched donor on Alpaca, SQuAD v2, and GSM8K, with an additional column showing the average across the three pools.

\begingroup
\small
\setlength{\tabcolsep}{1.8pt}
\renewcommand{\arraystretch}{1.05}
\begin{longtable}{p{0.3\textwidth}rrrrrrrrr}
\caption{SER results on differentiable transplant operators. Rows are grouped by the underlying shared-basis operator.}\label{tab:app:altops:ser} \\
\toprule
Pair (base$\leftarrow$donor) & $\lambda$ & \multicolumn{2}{c}{Alpaca} & \multicolumn{2}{c}{SQuAD v2} & \multicolumn{2}{c}{GSM8K} & \multicolumn{2}{c}{Avg} \\
 & & Base & Donor & Base & Donor & Base & Donor & Base & Donor \\
\cmidrule(lr){3-4}\cmidrule(lr){5-6}\cmidrule(lr){7-8}\cmidrule(lr){9-10}
\midrule
\endfirsthead
\toprule
Pair (base$\leftarrow$donor) & $\lambda$ & \multicolumn{2}{c}{Alpaca} & \multicolumn{2}{c}{SQuAD v2} & \multicolumn{2}{c}{GSM8K} & \multicolumn{2}{c}{Avg} \\
 & & Base & Donor & Base & Donor & Base & Donor & Base & Donor \\
\cmidrule(lr){3-4}\cmidrule(lr){5-6}\cmidrule(lr){7-8}\cmidrule(lr){9-10}
\midrule
\endhead
\midrule \multicolumn{10}{r}{\emph{Continued on next page}} \\
\endfoot
\bottomrule
\endlastfoot
\multicolumn{10}{l}{\textbf{FOCUS}} \\
\addlinespace[2pt]
Q3-0.6B $\leftarrow$ Gem2-2B & 16 & .6758 & .0000 & .7539 & .0000 & .9336 & .0000 & .7878 & .0000 \\
Q3-0.6B $\leftarrow$ Gem2-2B & 8 & .6680 & .0000 & .7578 & .0000 & .9141 & .0000 & .7799 & .0000 \\
Q3-0.6B $\leftarrow$ Gem2-2B & 4 & .6484 & .0000 & .7148 & .0000 & .9180 & .0000 & .7604 & .0000 \\
Q3-0.6B $\leftarrow$ Gem2-2B & 32 & .5938 & .0000 & .7305 & .0000 & .8984 & .0000 & .7409 & .0000 \\
Q3-0.6B $\leftarrow$ Gem2-2B & 64 & .5898 & .0000 & .6641 & .0000 & .7539 & .0000 & .6693 & .0000 \\
L3.2-1B $\leftarrow$ Q3-0.6B & 8 & .3555 & .0117 & .3477 & .0078 & .9453 & .0195 & .5495 & .0130 \\
Q3-0.6B $\leftarrow$ Gem2-2B & 128 & .4844 & .0000 & .5508 & .0000 & .5039 & .0000 & .5130 & .0000 \\
Q3-0.6B $\leftarrow$ Q2-0.5B & 4 & .2070 & .0039 & .2188 & .0039 & .9375 & .0352 & .4544 & .0143 \\
Q3-0.6B $\leftarrow$ Q2-0.5B & 8 & .2031 & .0000 & .2344 & .0195 & .9336 & .0430 & .4570 & .0208 \\
Q3-0.6B $\leftarrow$ Q2-0.5B & 2 & .2148 & .0039 & .2344 & .0000 & .9336 & .0820 & .4609 & .0286 \\
Q3-0.6B $\leftarrow$ Q2-0.5B & 16 & .1953 & .0000 & .2344 & .0195 & .9336 & .1094 & .4544 & .0430 \\
Q2-0.5B $\leftarrow$ Q3-0.6B & 16 & .3594 & .0508 & .2891 & .0469 & .9453 & .4023 & .5312 & .1667 \\
L3.2-1B $\leftarrow$ Gem2-2B & 2 & .3477 & .0000 & .4258 & .0000 & .3008 & .0000 & .3581 & .0000 \\
Q2-0.5B $\leftarrow$ Q3-0.6B & 32 & .3359 & .0117 & .2695 & .0508 & .9258 & .4102 & .5104 & .1576 \\
L3.2-1B $\leftarrow$ Gem2-2B & 4 & .2656 & .0000 & .3320 & .0000 & .1836 & .0000 & .2604 & .0000 \\
Q2-0.5B $\leftarrow$ Gem2-2B & 16 & .2109 & .0000 & .3086 & .0000 & .0391 & .0000 & .1862 & .0000 \\
\midrule
\multicolumn{10}{l}{\textbf{CLP}} \\
\addlinespace[2pt]
Q2-0.5B $\leftarrow$ Gem2-2B & 2 & .8516 & .0000 & .9648 & .0000 & .9922 & .0000 & .9362 & .0000 \\
Q3-0.6B $\leftarrow$ Gem2-2B & 2 & .8516 & .0000 & .9062 & .0000 & .9922 & .0000 & .9167 & .0000 \\
L3.2-1B $\leftarrow$ Q3-0.6B & 32 & .4570 & .0586 & .5273 & .0273 & .9805 & .0000 & .6549 & .0286 \\
L3.2-1B $\leftarrow$ Q3-0.6B & 16 & .4219 & .0352 & .5039 & .0156 & .9805 & .0000 & .6354 & .0169 \\
Gem2-2B $\leftarrow$ L3.2-1B & 2 & .5898 & .0000 & .2734 & .0000 & .4805 & .0000 & .4479 & .0000 \\
Gem2-2B $\leftarrow$ L3.2-1B & 4 & .5742 & .0000 & .3203 & .0000 & .4453 & .0000 & .4466 & .0000 \\
Q2-0.5B $\leftarrow$ Gem2-2B & 64 & .1328 & .0000 & .0977 & .0000 & .0234 & .0000 & .0846 & .0000 \\
L3.2-1B $\leftarrow$ Gem2-2B & 64 & .1445 & .0000 & .0391 & .0000 & .0547 & .0000 & .0794 & .0000 \\
L3.2-1B $\leftarrow$ Gem2-2B & 128 & .1406 & .0000 & .0430 & .0000 & .0508 & .0000 & .0781 & .0000 \\
L3.2-1B $\leftarrow$ Gem2-2B & 32 & .1289 & .0000 & .0391 & .0000 & .0430 & .0000 & .0703 & .0000 \\
Q2-0.5B $\leftarrow$ Gem2-2B & 128 & .0898 & .0000 & .0742 & .0000 & .0234 & .0000 & .0625 & .0000 \\
\midrule
\multicolumn{10}{l}{\textbf{WECHSEL}} \\
\addlinespace[2pt]
Q2-0.5B $\leftarrow$ Gem2-2B & 4 & .5078 & .0000 & .9336 & .0000 & .6914 & .0000 & .7109 & .0000 \\
Q2-0.5B $\leftarrow$ Gem2-2B & 8 & .3320 & .0000 & .7070 & .0000 & .2695 & .0000 & .4362 & .0000 \\
Q3-0.6B $\leftarrow$ L3.2-1B & 8 & .1016 & .0000 & .1367 & .0000 & .6055 & .0000 & .2812 & .0000 \\
Q3-0.6B $\leftarrow$ L3.2-1B & 2 & .1055 & .0000 & .1094 & .0000 & .5781 & .0000 & .2643 & .0000 \\
Q3-0.6B $\leftarrow$ L3.2-1B & 4 & .0469 & .0000 & .0859 & .0000 & .5156 & .0000 & .2161 & .0000 \\
Q3-0.6B $\leftarrow$ L3.2-1B & 16 & .0586 & .0000 & .0664 & .0000 & .2930 & .0000 & .1393 & .0000 \\
Q2-0.5B $\leftarrow$ Gem2-2B & 32 & .1328 & .0000 & .2148 & .0000 & .0156 & .0000 & .1211 & .0000 \\
Q2-0.5B $\leftarrow$ L3.2-1B & 4 & .1484 & .0000 & .0742 & .0000 & .0781 & .0000 & .1003 & .0000 \\
Q2-0.5B $\leftarrow$ L3.2-1B & 2 & .1250 & .0000 & .0781 & .0000 & .0938 & .0000 & .0990 & .0000 \\
Q2-0.5B $\leftarrow$ L3.2-1B & 8 & .1289 & .0000 & .0859 & .0000 & .0742 & .0000 & .0964 & .0000 \\
L3.2-1B $\leftarrow$ Gem2-2B & 8 & .0820 & .0000 & .1172 & .0000 & .0820 & .0000 & .0938 & .0000 \\
L3.2-1B $\leftarrow$ Gem2-2B & 16 & .0664 & .0000 & .0898 & .0000 & .0938 & .0000 & .0833 & .0000 \\
Q2-0.5B $\leftarrow$ L3.2-1B & 16 & .0977 & .0000 & .0586 & .0000 & .0469 & .0000 & .0677 & .0000 \\
\end{longtable}
\endgroup


\noindent\textbf{Analysis.}
Table~\ref{tab:app:altops:ser} shows that the same \emph{asymmetric realizability} pattern appears across differentiable shared-basis variants: the attacked base exhibits substantial SER while the patched donor remains near-silent.
The strongest examples occur with Gem2-2B as the donor, where both CLP and FOCUS yield high base SER across pools with near-zero donor SER.
WECHSEL achieves zero donor SER on every listed row; its base activation is smaller than CLP/FOCUS but remains substantial on multiple pairs.
Per-task variation reflects operator geometry: on closely related pairs, FOCUS shows nonzero donor SER on individual pools while keeping high base SER, while CLP exhibits the cleanest activation--stealth gap on these rows.
The donor-suppression objective transfers across shared-basis variants; the precise activation--stealth trade-off depends on the transplant operator's geometry.
